% GENERATED FILE. Edit canonical lessons, then run npm run generate:books.
% canonical-source-hash: fnv1a64:88731a424df38593
% canonical-lessons: JA-C167-hajimaru, JA-C167-hajimeru, JA-C167-oeru, JA-C167-kawaru, JA-C167-chuushi

\chapter{Starting, Finishing and Calling Off}
\label{ch:ja-starting-finishing-and-calling-off}
\hlchaptermodality{167}

\begin{chapteropening}
\textbf{By the end of this chapter:} \emph{I can say what will start, finish, change or be called off.}

Complete the last lesson of chapter 167: I can say what will start, finish, change or be called off.
\end{chapteropening}

\section[\texorpdfstring{\ja{はじまる} (hajimaru)}{はじまる (hajimaru)}]{\ja{はじまる} (hajimaru) — to begin}

\label{lesson:JA-C167-hajimaru}



\begin{quote}
Before the new one: say the Japanese for this evening, then the Japanese for every night.
\end{quote}

\subsection*{You'll want to know: \ja{はじまる}}
\textbf{\ja{はじまる}} — \emph{hajimaru} — \textquotedblleft{}to begin\textquotedblright{}.

\begin{quote}
\textbf{\ja{あした} \ja{まつりが} \ja{はじまります。}} — \emph{ashita matsuri ga hajimarimasu} — \textquotedblleft{}The festival begins tomorrow.\textquotedblright{}
\end{quote}

Polite \textbf{\ja{はじまります}}, \emph{hajimarimasu}; a godan verb.

\begin{cousinweb}
\textbf{\ja{はじまる}} (\emph{hajimaru}) is a native Japanese word (wago), inherited rather than borrowed from Chinese; its kanji was attached later as a kun'yomi, and its older origin is not settled (JMdict, KANJIDIC2).
\end{cousinweb}

\begin{grammarlens}[title={What you've built}]
One more word for everyday life at A2.
\end{grammarlens}

\subsection*{Your turn}
\begin{itemize}
\raggedright
  \item \emph{Say it:} \emph{hajimaru}
  \item \emph{Say it:} \emph{hajimaru}, once more
  \item \emph{Say it:} \emph{ashita matsuri ga hajimarimasu} — The festival begins tomorrow.
  \item \emph{Recall:} say the Japanese for this evening, then the Japanese for every night, then say \emph{hajimaru} again
\end{itemize}

\begin{tcolorbox}[breakable,colback=teal!4,colframe=teal!35!black,title={Before you move on}]
What does \ja{はじまる} mean? (\textquotedblleft{}To begin\textquotedblright{}.) Say it once more.
\end{tcolorbox}
\section[\texorpdfstring{\ja{はじめる} (hajimeru)}{はじめる (hajimeru)}]{\ja{はじめる} (hajimeru) — to start (something)}

\label{lesson:JA-C167-hajimeru}



\begin{quote}
Before the new one: say the Japanese for every night, then the Japanese for to begin.
\end{quote}

\subsection*{You'll want to know: \ja{はじめる}}
\textbf{\ja{はじめる}} — \emph{hajimeru} — \textquotedblleft{}to start (something)\textquotedblright{}.

\begin{quote}
\textbf{\ja{らいしゅう} \ja{しごとを} \ja{はじめます。}} — \emph{raishū shigoto o hajimemasu} — \textquotedblleft{}I will start work next week.\textquotedblright{}
\end{quote}

Polite \textbf{\ja{はじめます}}, \emph{hajimemasu}; an ichidan verb. Its intransitive partner is \textbf{\ja{はじまる}}, \emph{hajimaru}.

\begin{cousinweb}
\textbf{\ja{はじめる}} (\emph{hajimeru}) is a native Japanese word (wago), inherited rather than borrowed from Chinese, and built on the same root as \emph{hajimaru}, \textquotedblleft{}to begin\textquotedblright{}; its kanji was attached later as a kun'yomi (JMdict, KANJIDIC2).
\end{cousinweb}

\begin{grammarlens}[title={What you've built}]
One more word for everyday life at A2.
\end{grammarlens}

\subsection*{Your turn}
\begin{itemize}
\raggedright
  \item \emph{Say it:} \emph{hajimeru}
  \item \emph{Say it:} \emph{hajimeru}, once more
  \item \emph{Say it:} \emph{raishū shigoto o hajimemasu} — I will start work next week.
  \item \emph{Recall:} say the Japanese for every night, then the Japanese for to begin, then say \emph{hajimeru} again
\end{itemize}

\begin{tcolorbox}[breakable,colback=teal!4,colframe=teal!35!black,title={Before you move on}]
What does \ja{はじめる} mean? (\textquotedblleft{}To start (something)\textquotedblright{}.) Say it once more.
\end{tcolorbox}
\section[\texorpdfstring{\ja{おえる} (oeru)}{おえる (oeru)}]{\ja{おえる} (oeru) — to finish (something)}

\label{lesson:JA-C167-oeru}



\begin{quote}
Before the new one: say the Japanese for to begin, then the Japanese for to start.
\end{quote}

\subsection*{You'll want to know: \ja{おえる}}
\textbf{\ja{おえる}} — \emph{oeru} — \textquotedblleft{}to finish (something)\textquotedblright{}.

\begin{quote}
\textbf{\ja{あした} \ja{しごとを} \ja{おえます。}} — \emph{ashita shigoto o oemasu} — \textquotedblleft{}I will finish the work tomorrow.\textquotedblright{}
\end{quote}

Polite \textbf{\ja{おえます}}, \emph{oemasu}; an ichidan verb.

\begin{cousinweb}
\textbf{\ja{おえる}} (\emph{oeru}) is a native Japanese word (wago), inherited rather than borrowed from Chinese, and built on the same root as \emph{owaru}, \textquotedblleft{}to end\textquotedblright{}; its kanji was attached later as a kun'yomi (JMdict, KANJIDIC2).
\end{cousinweb}

\begin{grammarlens}[title={What you've built}]
One more word for everyday life at A2.
\end{grammarlens}

\subsection*{Your turn}
\begin{itemize}
\raggedright
  \item \emph{Say it:} \emph{oeru}
  \item \emph{Say it:} \emph{oeru}, once more
  \item \emph{Say it:} \emph{ashita shigoto o oemasu} — I will finish the work tomorrow.
  \item \emph{Recall:} say the Japanese for to begin, then the Japanese for to start, then say \emph{oeru} again
\end{itemize}

\begin{tcolorbox}[breakable,colback=teal!4,colframe=teal!35!black,title={Before you move on}]
What does \ja{おえる} mean? (\textquotedblleft{}To finish (something)\textquotedblright{}.) Say it once more.
\end{tcolorbox}
\section[\texorpdfstring{\ja{かわる} (kawaru)}{かわる (kawaru)}]{\ja{かわる} (kawaru) — to become different; to change}

\label{lesson:JA-C167-kawaru}



\begin{quote}
Before the new one: say the Japanese for to start, then the Japanese for to finish.
\end{quote}

\subsection*{You'll want to know: \ja{かわる}}
\textbf{\ja{かわる}} — \emph{kawaru} — \textquotedblleft{}to become different; to change\textquotedblright{}.

\begin{quote}
\textbf{\ja{らいげつ} \ja{しごとの} \ja{じかんが} \ja{かわります。}} — \emph{raigetsu shigoto no jikan ga kawarimasu} — \textquotedblleft{}Next month my working hours will change.\textquotedblright{}
\end{quote}

Polite \textbf{\ja{かわります}}, \emph{kawarimasu}; a godan verb.

\begin{cousinweb}
\textbf{\ja{かわる}} (\emph{kawaru}) is a native Japanese word (wago), inherited rather than borrowed from Chinese; its kanji was attached later as a kun'yomi, and its older origin is not settled (JMdict, KANJIDIC2).
\end{cousinweb}

\begin{grammarlens}[title={What you've built}]
One more word for everyday life at A2.
\end{grammarlens}

\subsection*{Your turn}
\begin{itemize}
\raggedright
  \item \emph{Say it:} \emph{kawaru}
  \item \emph{Say it:} \emph{kawaru}, once more
  \item \emph{Say it:} \emph{raigetsu shigoto no jikan ga kawarimasu} — Next month my working hours will change.
  \item \emph{Recall:} say the Japanese for to start, then the Japanese for to finish, then say \emph{kawaru} again
\end{itemize}

\begin{tcolorbox}[breakable,colback=teal!4,colframe=teal!35!black,title={Before you move on}]
What does \ja{かわる} mean? (\textquotedblleft{}To become different; to change\textquotedblright{}.) Say it once more.
\end{tcolorbox}
\section[\texorpdfstring{\ja{ちゅうし} (chūshi)}{ちゅうし (chūshi)}]{\ja{ちゅうし} (chūshi) — calling off; suspension}

\label{lesson:JA-C167-chuushi}



\begin{quote}
Before the new one: say the Japanese for to finish, then the Japanese for to become different.
\end{quote}

\subsection*{You'll want to know: \ja{ちゅうし}}
\textbf{\ja{ちゅうし}} — \emph{chūshi} — \textquotedblleft{}calling off; suspension\textquotedblright{}.

\begin{quote}
\textbf{\ja{あしたの} \ja{まつりは} \ja{ちゅうしでしょう。}} — \emph{ashita no matsuri wa chūshi deshō} — \textquotedblleft{}Tomorrow's festival will probably be called off.\textquotedblright{}
\end{quote}

\begin{cousinweb}
\textbf{\ja{ちゅうし}} (\emph{chūshi}) is Sino-Japanese (kango): two Chinese characters read with their on'yomi, readings borrowed from Chinese. They are \emph{chū}, \textquotedblleft{}middle\textquotedblright{} (Middle Chinese \emph{trjuwng}) and \emph{shi}, \textquotedblleft{}stop\textquotedblright{} (\emph{tsyiX}) (KANJIDIC2; Middle Chinese in Baxter's transcription).
\end{cousinweb}

\begin{grammarlens}[title={What you've built}]
One more word for everyday life at A2.
\end{grammarlens}

\subsection*{Your turn}
\begin{itemize}
\raggedright
  \item \emph{Say it:} \emph{chūshi}
  \item \emph{Say it:} \emph{chūshi}, once more
  \item \emph{Say it:} \emph{ashita no matsuri wa chūshi deshō} — Tomorrow's festival will probably be called off.
  \item \emph{Recall:} say the Japanese for to finish, then the Japanese for to become different, then say \emph{chūshi} again
\end{itemize}

\begin{tcolorbox}[breakable,colback=teal!4,colframe=teal!35!black,title={Before you move on}]
What does \ja{ちゅうし} mean? (\textquotedblleft{}Calling off; suspension\textquotedblright{}.) Say it once more.
\end{tcolorbox}
